% ===================================================================
%  TABELLA nr. 4 — Età madira e vegliadetgna.
%  Traduction 2026 d'apres texto/io, controlee sur texto/fr, texto/ca
%  et texto/oc. Consigne : voir texto/rm/10-tabella-01.tex.
%
%  DEUX TABLEAUX PLUS TOT, LE ROMANCHE REFUSAIT TESTA POUR LA TETE ET
%  GARDAIT CAPUT ; SUR CETTE PAGE-CI IL FAIT L'INVERSE, ET C'EST LE
%  MEME MOT DE FAMILLE QUI LE MONTRE DEUX FOIS.
%
%  Les mots de parente que cette page emploie se rangent d'eux-memes
%  en deux tas, et le partage n'est pas celui qu'on attendait.
%
%  Nourrice, tous nes de la bouche d'un enfant :
%    « bab »    le pere      — PATER a disparu ;
%    « mamma »  la mere      — MATER a disparu ;
%    « tat »    le grand-pere ;
%    « tatta »  la grand-mere ;
%    « bistat » l'arriere-grand-pere, qui est « bis » + « tat » —
%               le mot savant lui-meme est bati sur le mot d'enfant.
%
%  Latins, et souvent plus vieux que chez les voisins :
%    « schuègr », « schuegra » de SOCER et SOCRA, le beau-pere et la
%       belle-mere, la ou le francais a fabrique « beau-pere » ;
%    « quinà », « quinada » de COGNATUS ;
%    « niera » de NURUS ;
%    « abiadi », « abiadia » le petit-fils et la petite-fille, de
%       AVIATICUS — un mot que ni l'italien ni le francais n'ont
%       garde, l'un disant « nipote » et l'autre un compose.
%
%  LA LIGNE NE PASSE PAS ENTRE LE PROCHE ET LE LOINTAIN, NI ENTRE
%  L'ANCIEN ET LE RECENT : ELLE PASSE ENTRE LES MOTS DONT ON APPELLE
%  QUELQU'UN ET LES MOTS DONT ON PARLE DE QUELQU'UN. On crie « bab »,
%  « mamma », « tat », « tatta » ; personne n'a jamais appele
%  quelqu'un « abiadi », « quinà » ou « niera » — ce sont des mots
%  qu'on emploie A PROPOS d'une personne, jamais A une personne. Le
%  vocatif est de nourrice, la reference est latine.
%
%  ET CELA REPOND AU TABLEAU 2. Le romanche n'est pas « conservateur »
%  en general : il a garde CAPUT parce que « chau » est un mot dont on
%  parle, et il a perdu PATER parce que « bab » est un mot dont on
%  appelle. La conservation et l'innovation ne se partagent pas les
%  domaines, elles se partagent LES EMPLOIS. Le tableau 16
%  luxembourgeois avait trouve la ligne entre ce qui a un guichet et
%  ce qui n'en a pas ; ici elle passe entre l'adresse et la mention,
%  et c'est plus fin d'un cran.
%
%  « bandoliera vildo-sako » de l'alinea 4 est mis TOUT ENTIER en
%  gras, comme dans les vingt-et-une colonnes precedentes : le
%  fac-simile ido n'en grasse que la seconde moitie.
% ===================================================================

%%K t04-apar-1 apar
\VUsaut{4.0mm}
\VUtitre{15.0pt}{60}{TABELLA nr. 4}
\VUsaut{1.6mm}
\VUtitre{11.4pt}{0}{Età madira e vegliadetgna.}
\VUsaut{3.2mm}

%%K t04-tit-1 sub
\VUcentre{11.4pt}{0}{\textit{Emprima scena.}}
\VUsaut{1.6mm}
\VUtitre{11.4pt}{0}{Il banchet da nozzas.}
\VUsaut{2.6mm}

%%K t04-tit-2 sub
\VUpk{10.2pt}{40}{L'atun.}
\VUsaut{3.0mm}

%%K t04-c1-01-1 p
1. --- Ils giuvens èn daventads pli gronds. In d'els, Gion, sa marida. Nus
assistain al \VUgras{banchet da nozzas}, che raduna, enturn la medema maisa, las
famiglias dals spus.

%%K t04-c1-02-1 p
2. --- Nus essan en \VUgras{atun}, en il \VUgras{mais da settember}. Il vent fraid
d'\VUgras{october} e da \VUgras{november} n'ha anc betg spoglià la champagna da ses
fegliom verd. L'aria da la saira è tievda e parfumada; perquai ha la tschaina lieu
sut la \VUgras{veranda} \textsuperscript{(8)} en l'iert.

%%K t04-c1-03-1 p
3. --- Pli lunsch, sin la \VUgras{collina} \textsuperscript{(3)}, sa chatta la chasa
dals geniturs da Gion, in elegant \VUgras{chastè} \textsuperscript{(1)} zuppà en il verd;
bellas vignas, che dattan in vin excellent, cuvran il pendent da la collina. Il
vignerun als culvitescha ed als tgira.

%%K t04-c1-04-1 p
4. --- Sur las vignas passa in svol da \VUgras{pernischs} \textsuperscript{(2)}
spaventads da l'avischinanza da la \VUgras{guardia dal chatsch} \textsuperscript{(5)}.
Quel, cun il \VUgras{sclop} sut il bratsch e cun in \VUgras{satg da chatscha en
bandulera}, explorescha la \VUgras{buschaglia} \textsuperscript{(4)} cun ses
\VUgras{chaun} \textsuperscript{(6)}. El survegliescha la tenuta ed impedescha ils
chatschaders da furtim da destruir la selvaschina.

%%K t04-c1-05-1 p
5. --- Ordaifer la \VUgras{veranda} rampigna in \VUgras{spalier da vit}, dal qual
pendan densas \VUgras{rapas} \textsuperscript{(7)}.

%%K t04-c1-06-1 p
6. --- Las bellas flurs d'atun, che ornan la \VUgras{maisa} \textsuperscript{(16)}, èn
\VUgras{crisantems} \textsuperscript{(9)}; il \VUgras{bab} \textsuperscript{(10)} da la
spusa ha mess ina d'ellas en il fori dal buttun da ses frac.

%%K t04-c1-07-1 p
7. --- Il past va a fin. Il \VUgras{servitur} \textsuperscript{(11)} cumenza a purtar
ils plats da dessert e vegn prest a prender davent las differentas parts da
l'ustensiglia da maisa, \VUgras{chadagnas} \textsuperscript{(14)}, \VUgras{furtgettas}
\textsuperscript{(12)}, \VUgras{curtels} \textsuperscript{(13)}, \VUgras{plats gronds}
\textsuperscript{(15)}, ch'ins n'ha betg pli da basegn.

%%K t04-tit-3 sub
\VUpk{10.2pt}{40}{La parentella.}
\VUsaut{3.0mm}

%%K t04-c1-08-1 p
8. --- Tuts paran allegers, e il surrir cun il qual la \VUgras{mamma}
\textsuperscript{(17)} dal spus guarda ses uffants cumprova sia fortuna.

%%K t04-c1-09-1 p
9. --- Quellas nozzas uneschan duas famiglias ritgas dal pajais. Gion, il
\VUgras{spus} \textsuperscript{(18)}, è \VUgras{figl} d'in grond possessur ed el daventa
il \VUgras{genderin} d'in industrial abil.

%%K t04-c1-10-1 p
10. --- Ils \VUgras{spus} èn giuvens e tuttina eran els \VUgras{spusads} dapi ditg.
La \VUgras{giuvna dunna} \textsuperscript{(19)}, Marta, è charmanta en sia
\VUgras{rassa alva} ornada cun flurs d'oranger.

%%K t04-c1-11-1 p
11. --- Ses \VUgras{schuègr} \textsuperscript{(20)} è fitg fortunà d'avair ina
\VUgras{niera} uschè gentila, e la \VUgras{schuegra} \textsuperscript{(21)} da Gion
surri vesend uschè bella sia \VUgras{figlia}.

%%K t04-c1-12-1 p
12. --- A la fin da la maisa sesan ils auters \VUgras{envidads}
\textsuperscript{(22)}, \VUgras{parents, amis, cusrins, cusrinas}. In
\VUgras{maistral da maisa} \textsuperscript{(23)}, cun la servietta sut il bratsch,
als serva diligentamain.

%%K t04-tit-4 sub
\VUpk{10.2pt}{40}{Ils amis.}
\VUsaut{3.0mm}

%%K t04-c1-13-1 p
13. --- Da la vart dretga da la maisa, qua è ina \VUgras{damaisella}
\textsuperscript{(24)}, \VUgras{amia} da la giuvna spusa, lura ses \VUgras{frar}, ch'è
ses \VUgras{cumpogn d'onur} \textsuperscript{(25)}. El tschintscha cun sia
\VUgras{damaisella d'onur} \textsuperscript{(26)}, sora dal spus, che daventa cun
quellas nozzas la \VUgras{quinada} da Marta. Ella è ina charmanta giuvna. Ella ha
bels \VUgras{gioghels}, ina \VUgras{culiera da perlas} ed in \VUgras{bratschadi}
d'aur.

%%K t04-c1-14-1 p
14. --- Sper els è il signur ch'è en pes ed auza ses magiel in \VUgras{ami}
\textsuperscript{(27)} da la famiglia. El è in dals \VUgras{perditgs dal spus}. El
\VUgras{fa in brindis}, el beva a la \VUgras{sanadad} dals giuvens spus ed als
giavischa ina lunga fortuna. El è vestgì cun il costum da ceremonia, cun in frac e
cun \VUgras{chalzers lustrads} \textsuperscript{(28)}. El accumpogna la \VUgras{tatta}
\textsuperscript{(29)}, ch'è \VUgras{vaidgua}.

%%K t04-c1-15-1 p
15. --- Il giuvenot en \VUgras{frac} \textsuperscript{(31)}, cun mustatschas nairas e
sutgladas, è il \VUgras{quinà} \textsuperscript{(30)} da Gion. El è in inschigner
eminent. El è il vischin d'ina \VUgras{giuvna dama} \textsuperscript{(33)} fitg
eleganta, la \VUgras{sora pli veglia} da Marta e la mamma da la charmanta
mattinetta, che vegn, dond il bratsch a ses cusrin, ad offrir ina flur ed a
\VUgras{giavischar} a el \VUgras{la buna saira}. Quella dama ha fitg bellas
\VUgras{urengias} ed in elegant \VUgras{chapè}.

%%K t04-c1-16-1 p
16. --- Il pitschen cusrin, uschè curius cun sia \VUgras{blusa} \textsuperscript{(36)}
e sias sfunsadas \VUgras{chautschas curtas} \textsuperscript{(37)}, fa fitg
cumpassiun. El è \VUgras{orfen} \textsuperscript{(35)}. Fitg giuven ha el pers ses
bab e sia mamma. Ses aug è ses \VUgras{tutur} \textsuperscript{(38)}, ed el è restà
nunmaridà per educar il pauprin. El è in um bunmanaivel. El è in pau chalv e fitg
curt da vista; el duvra \VUgras{binoclas dal nas}.

%%K t04-tit-5 sub
\VUsaut{2.4mm}
\VUpk{10.2pt}{40}{Il dessert.}
\VUsaut{3.0mm}

%%K t04-c1-17-1 p
17. --- Davos ils envidads, sin ina maisa cuvrida cun ina \VUgras{tavatscha}, è
preparà il \VUgras{dessert} \textsuperscript{(39)}. En ina gronda cuppa, qua èn
fritgas, \VUgras{persis} \textsuperscript{(40)}, \VUgras{pirs} \textsuperscript{(41)},
\VUgras{mails} \textsuperscript{(42)}, \VUgras{abricos} \textsuperscript{(43)} e
\VUgras{rapas} \textsuperscript{(44)}. Dasperas, \VUgras{ananas}
\textsuperscript{(45)}, ina bella \VUgras{tuorta} \textsuperscript{(46)},
\VUgras{buttiglias da liquor} \textsuperscript{(48)} e \VUgras{schampagner}
\textsuperscript{(49)} ed era colonnas da \VUgras{plats} \textsuperscript{(47)} nets.

%%K t04-c1-18-1 p
18. --- Il \VUgras{sumelier} \textsuperscript{(50)} stapa ina \VUgras{buttiglia}
\textsuperscript{(52)} da vin vegl cun in \VUgras{tiratap} \textsuperscript{(51)}. Il
\VUgras{tap da suver} \textsuperscript{(53)} para fitg dur da trair. Quel sumelier ha
ina \VUgras{servietta} \textsuperscript{(54)} sut il bratsch.

%%K t04-c1-19-1 p
19. --- Suenter la tschaina van ils signurs en la chombra da fimar, e las dunnas
van a sa preparar per il bal.

%%K t04-tit-6 sub
\VUsaut{5.0mm}
\VUcentre{11.4pt}{0}{\textit{Segunda scena.}}
\VUsaut{1.6mm}
\VUpk{11.4pt}{40}{L'anniversari dal tat.}
\VUsaut{2.6mm}

%%K t04-tit-7 sub
\VUpk{10.2pt}{40}{La visita.}
\VUsaut{3.0mm}

%%K t04-c2-01-1 p
1. --- Diesch onns èn passads, ils geniturs da Gion èn daventads vegls e han fitg
gugent lur trais abiadis, dus \VUgras{abiadis} \textsuperscript{(2)} ed ina
\VUgras{abiadia} \textsuperscript{(3)}. Perquai ch'oz è \VUgras{l'anniversari dal
tat}, vegnan ils uffants a lur giavischar in bun di da naschientscha.

%%K t04-c2-02-1 p
2. --- Il pitschen Peider è anc fitg giuven, el ha mo trais onns. El vesa il
\VUgras{giat} \textsuperscript{(1)} sa avischinar. El vul charezzar ses bel
\VUgras{pel} da seida e sia lunga \VUgras{cua}; el na patratga betg che sut ils
graziusa \VUgras{peis} sa zuppan \VUgras{unglas} pizzas e malignas, che al vegnan
forsa a sgrafar.

%%K t04-c2-03-1 p
3. --- Sia sora Gabriela, che al dat il maun, è bler pli seriusa, e Marcel cun ses
grond \VUgras{mantè} \textsuperscript{(4)} e sia \VUgras{tschinta}
\textsuperscript{(5)} da corom para in pau daventà um, malgrà sias chautschas curtas,
che laschan vesair sias \VUgras{chaltschas} \textsuperscript{(6)}.

%%K t04-c2-04-1 p
4. --- Lur bab ha mess ses gross \VUgras{surtut} \textsuperscript{(7)} d'enviern cun
\VUgras{collar da pelitscha} \textsuperscript{(8)} e, en sia \VUgras{giaglioffa}
\textsuperscript{(9)}, ha el in foulard. Sia dunna ha ses chapè da feltr ornà cun
\VUgras{plimas} \textsuperscript{(10)}, sia \VUgras{giacca} \textsuperscript{(11)}, sia
\VUgras{boa da pelitscha} \textsuperscript{(12)} e ses \VUgras{manschun}
\textsuperscript{(13)}.

%%K t04-c2-05-1 p
5. --- I fa fitg fraid, perquai ch'i regia l'enviern. La \VUgras{naiv}
\textsuperscript{(19)} è crudada tut il di ed ils tetgs da las chasas èn tuts alvs.
Cun la \VUgras{notg} daventa il fraid anc pli asper. En il tschiel splenduran la
\VUgras{glina} \textsuperscript{(17)} e las \VUgras{stailas} \textsuperscript{(18)} e
traversan la \VUgras{stgiradetgna}.

%%K t04-tit-8 sub
\VUsaut{4.2mm}
\VUpk{10.2pt}{40}{La malsogna.}
\VUsaut{3.0mm}

%%K t04-c2-06-1 p
6. --- Il paupra \VUgras{tschorv} \textsuperscript{(20)}, che traversa la plazza avant
l'\VUgras{apoteca} \textsuperscript{(16)}, manà da ses \VUgras{pudel}
\textsuperscript{(21)}, è fitg malfortunà. Justina, la \VUgras{serventa}
\textsuperscript{(14)}, porta or da la cuschina ina \VUgras{scadiola}
\textsuperscript{(15)} da \VUgras{tisana} fitg chauda e zuccrada generusamain.

%%K t04-c2-07-1 p
7. --- La \VUgras{tatta} \textsuperscript{(23)}, che vul tschertgar ses
\VUgras{lorgnon} \textsuperscript{(26)} sin la maisa per leger la \VUgras{recepta}
\textsuperscript{(27)} ch'il \VUgras{medi} \textsuperscript{(25)} ha gist scrit, sa
auza da ses sutgun, sa sustegnend sin ses \VUgras{bastun} \textsuperscript{(24)}.

%%K t04-c2-08-1 p
8. --- Savens patescha ella da \VUgras{malsognettas, vertigins, catars, mal dals
dents}, sias chommas èn flaivlas e ses egls n'èn betg pli fitg buns. Malgrà quai fa
ella gugent gieu cun ses vegl \VUgras{medi}, che vul adina prescriver ad ella ina
\VUgras{poziun} \textsuperscript{(28)}. Oz è ella fitg allegra, perquai ch'ella ha
enturn sai sia giuvna famiglia.

%%K t04-c2-09-1 p
9. --- I è confortabel en la chombra bain serrada. En il \VUgras{chamin}
\textsuperscript{(29)} da marmel arda in \VUgras{bel fieu} \textsuperscript{(30)}, ed
ins sa senta fortunà en la dultscha \VUgras{glisch} da la \VUgras{lampa}
\textsuperscript{(31)} ch'ins ha gist emprendì.

%%K t04-tit-9 sub
\VUsaut{4.2mm}
\VUpk{10.2pt}{40}{Il tat.}
\VUsaut{3.0mm}

%%K t04-c2-10-1 p
10. --- Er il \VUgras{tat} \textsuperscript{(33)} ha emblidà ch'el era malsaun, che
ses \VUgras{rumatissem} al fixescha en ses \VUgras{sutgun} \textsuperscript{(34)} e
ch'el aveva anc ier \VUgras{fevra e svanids}. El na s'algorda betg pli ch'el ha
patì l'enviern passà d'ina \VUgras{pulmonitis} e ch'ina \VUgras{tuss} rauca e
penibla al ha fatg fitg patir.

%%K t04-c2-10-2 p
E tuttina è el mo cun grondas fadias vegnì sanà. Uss va el in pau meglier. Ma sia
chomma, ch'el posa sin in \VUgras{chaschin} \textsuperscript{(38)} mess sin in
pitschen \VUgras{stgabè} \textsuperscript{(39)}, al fa era mal.

%%K t04-c2-11-1 p
11. --- Cur ch'el è bain enviclà en sia chauda \VUgras{rassa da chombra}
\textsuperscript{(35)} cun gross \VUgras{buttuns} \textsuperscript{(36)} da
madraperla e cur ch'el ha sper sai, per al leger la gasetta, la pitschna
\VUgras{orfna} \textsuperscript{(40)} vestgida cun ina \VUgras{cuffa} alva, cun ina
blaua \VUgras{rassa} \textsuperscript{(42)} ed ina \VUgras{pelerina}
\textsuperscript{(41)}, na sa lamenta el betg.

%%K t04-c2-12-1 p
12. --- Mintg'onn, cur ch'il \VUgras{chalender} \textsuperscript{(43)} mussa danovamain
la \VUgras{data} dal emprim da \VUgras{december}, spetga el cun nunpazienza ses
\VUgras{abiadis}, perquai ch'el sa ch'els n'emblidan betg quella \VUgras{data}
impurtanta.

%%K t04-c2-13-1 p
13. --- El s'algorda cun emoziun dal temp fitg dalunsch, cur ch'el è ì persunalmain
a giavischar ina buna festa al gloerius \VUgras{mareschal} imperial, dal qual il
\VUgras{portret} \textsuperscript{(44)} penda vi da la paraid e ch'è il
\VUgras{bistat} da ses uffants.

\VUsaut{6.0mm}
\VUfilet{22mm}
